\chapter{Orbifolds}

\section{Quotients and basic structure}

\begin{example}[Reflection quotient]
\label{thurston-ch13-reflection-quotient}
\group{examples}
\source{thurston:page-0315,thurston:page-0316,thurston:page-0317,thurston:page-0318}
\notready
The quotient of Euclidean space by a properly discontinuous reflection group is
represented by a fundamental polyhedron with mirror faces.  The orientation-
preserving subgroup of the rectangular reflection group gives a spherical
orbifold with four order-two cone points, each of cone angle $\pi$; the
orientation-preserving half-turn group generated by the three indicated lattice
families gives $S^3$ with a Euclidean structure off a singular Borromean-ring
locus of cone angle $\pi$.
\end{example}
\begin{proof}
\sketch Fold a fundamental rectangle or cube along the reflection faces and
record the nontrivial isotropy.  The four rectangle corners yield cone angle
$\pi$; in the three-dimensional construction the axes descend to the three
linked singular circles.
\end{proof}

\begin{definition}[Orbifold atlas]
\label{thurston-ch13-orbifold-atlas}
\group{foundations}
\source{thurston:page-0318,thurston:page-0319}
\notready
An $n$-orbifold is a Hausdorff space $X_O$ covered by open sets $U_i$ closed
under finite intersections, with charts
$U_i\cong\widetilde U_i/\Gamma_i$, where $\widetilde U_i\subset\mathbb R^n$
is open and $\Gamma_i$ is finite.  For $U_i\subset U_j$ there are an injective
homomorphism $f_{ij}:\Gamma_i\hookrightarrow\Gamma_j$ and an equivariant
embedding $\widetilde\varphi_{ij}:\widetilde U_i\hookrightarrow\widetilde U_j$.
These data are defined up to the natural action of $\Gamma_j$; on triple
overlaps the two composites agree up to conjugacy.
\end{definition}
\begin{proof}
\sketch The quotient maps give the local models.  Restriction of a chart to a
smaller chart induces the injective isotropy map and equivariant lift; changing
the lift changes both maps by an element of the larger isotropy group.
\end{proof}

\begin{proposition}[Quotients are orbifolds]
\label{thurston-ch13-quotient-orbifold}
\dcref{13.2.1}
\group{foundations}
\source{thurston:page-0320}
\notready
If a group $\Gamma$ acts properly discontinuously on a manifold $M$, then
$M/\Gamma$ has a natural orbifold structure.  At the image of $\widetilde x$,
the chart group is the isotropy group of $\widetilde x$.
\end{proposition}
\begin{proof}
\sketch Choose a small isotropy-invariant neighborhood of each lift, disjoint
from its translates outside the isotropy group.  Finite intersections lift to
intersections of conjugate neighborhoods, with chart group the intersection of
the corresponding conjugate isotropy groups.
\end{proof}

\begin{definition}[Singular locus and covering orbifold]
\label{thurston-ch13-singular-covering}
\group{foundations}
\source{thurston:page-0320,thurston:page-0321}
\notready
The local isotropy group $\Gamma_x$ is well-defined up to isomorphism, and
$\Sigma_O=\{x:\Gamma_x\ne1\}$ is the singular locus.  A covering orbifold
$\widetilde O\to O$ is locally modeled over $\widetilde U/\Gamma$ by components
$\widetilde U/\Gamma_i$, with $\Gamma_i\leq\Gamma$, compatibly with projection.
\end{definition}
\begin{proof}
\sketch Common refinements identify local isotropy groups.  The singular set is
closed and nowhere dense because a nontrivial finite-order homeomorphism of a
manifold cannot fix an open set.  The stated local model is precisely the
covering condition.
\end{proof}

\begin{definition}[Good orbifold and fundamental group]
\label{thurston-ch13-good-orbifold}
\dcref{13.2.3,13.2.5}
\group{foundations}
\source{thurston:page-0322,thurston:page-0325}
\notready
An orbifold is good if it has a manifold covering orbifold, and bad otherwise.
Its fundamental group is the group of deck transformations of its universal
covering orbifold.  In particular, if $M$ is simply connected and $O=M/\Gamma$,
then $\pi_1(O)=\Gamma$.
\end{definition}
\begin{proof}
\sketch The definitions identify coverings with subgroups of the deck group.
The universal cover, when constructed below, has deck group acting transitively
on fibers, giving the asserted fundamental group.
\end{proof}

\begin{proposition}[Universal orbifold cover]
\label{thurston-ch13-universal-cover}
\dcref{13.2.4}
\group{foundations}
\source{thurston:page-0323,thurston:page-0324,thurston:page-0325}
\notready
Every orbifold $O$ has a connected universal covering orbifold
$\widetilde O\to O$: for a base point away from $\Sigma_O$, every pointed
covering orbifold over $O$ receives a unique pointed lift from $\widetilde O$.
\end{proposition}
\begin{proof}
\sketch Take the component of the base point in the fiber product of all
pointed connected covering orbifolds.  The resulting inverse-limit cover maps
to every member and has the lifting property.  In codimension at least two,
the same construction is equivalently obtained by completing a universal cover
of the nonsingular locus, with the local isotropy relations imposed.
\end{proof}

\begin{definition}[Orbifold boundary and suborbifold]
\label{thurston-ch13-boundary-suborbifold}
\dcref{13.2.6,13.2.7}
\group{foundations}
\source{thurston:page-0325,thurston:page-0326}
\notready
An orbifold with boundary is locally modeled on a half-space modulo a finite
group preserving the boundary hyperplane.  A suborbifold is a subspace whose
local inverse image in every chart is a submanifold invariant under the local
group, with the induced quotient charts.
\end{definition}
\begin{proof}
\sketch Restrict each local quotient chart to the half-space or invariant
submanifold.  Compatibility of the ambient charts supplies the induced atlas.
\end{proof}

\section{Two-dimensional orbifolds}

\begin{proposition}[Classification of two-dimensional singular strata]
\label{thurston-ch13-surface-singular-locus}
\dcref{13.3.1}
\group{surfaces}
\source{thurston:page-0326}
\notready
The singular locus of a two-orbifold is a union of mirror arcs and circles,
with isolated corner reflectors at endpoints of mirror arcs.  Away from mirrors,
isolated singular points have cyclic rotation groups.
\end{proposition}
\begin{proof}
\sketch Finite subgroups of the Euclidean isometry group fixing a point are
cyclic rotations or include reflections.  The fixed sets of reflections are
lines, and intersections of reflection lines give corner points.
\end{proof}

\begin{proposition}[Analytic orbifolds are good]
\label{thurston-ch13-analytic-good}
\dcref{13.3.2}
\group{surfaces}
\source{thurston:page-0328}
\notready
If $G$ is an analytic group of diffeomorphisms of a manifold $X$, every
$(G,X)$-orbifold is good and admits a developing map from its universal cover.
\end{proposition}
\begin{proof}
\sketch Analytic continuation of local lifts along paths produces a developing
map; analyticity makes continuation unique after isotropy identifications.
The resulting continuation cover is a manifold covering orbifold.
\end{proof}

\begin{definition}[Orbifold Euler characteristic]
\label{thurston-ch13-orbifold-euler}
\dcref{13.3.3}
\group{surfaces}
\source{thurston:page-0329}
\notready
For a cell division subordinate to the isotropy stratification, define
$\chi(O)=\sum_C(-1)^{\dim C}/|\Gamma_C|$, where $\Gamma_C$ is the isotropy
group on the interior of $C$.
\end{definition}
\begin{proof}
\sketch Lift the cell division to a manifold covering orbifold.  Each lifted
cell has $|\Gamma_C|$ translates over $C$, so ordinary Euler characteristic
divided by the covering degree gives the displayed weighted sum.
\end{proof}

\begin{proposition}[Euler characteristic under covers]
\label{thurston-ch13-euler-cover}
\dcref{13.3.4}
\group{surfaces}
\source{thurston:page-0329}
\notready
If $\widetilde O\to O$ is a $k$-sheeted covering orbifold, then
$\chi(\widetilde O)=k\chi(O)$.
\end{proposition}
\begin{proof}
\sketch Lift a subordinate cell division and count the $k$ copies of each
weighted cell, preserving isotropy weights.
\end{proof}

\begin{theorem}[Uniformization of closed two-orbifolds]
\label{thurston-ch13-two-orbifold-uniformization}
\dcref{13.3.6}
\group{surfaces}
\source{thurston:page-0330,thurston:page-0331,thurston:page-0332,thurston:page-0333,thurston:page-0334,thurston:page-0335}
\notready
A closed two-orbifold admits an elliptic, parabolic, or hyperbolic structure
if and only if it is good.  More precisely, it is hyperbolic exactly when
$\chi(O)<0$, it is parabolic exactly when $\chi(O)=0$, and it is either
elliptic or bad exactly when $\chi(O)>0$.
\end{theorem}
\begin{proof}
\sketch The weighted Gauss--Bonnet formula forces the asserted curvature sign.
Classifying the nonnegative-Euler-characteristic cases separates the bad
orbifolds from the explicit spherical and Euclidean quotients.  When the Euler
characteristic is negative, cutting into elementary orbifolds and gluing their
hyperbolic structures constructs the required metric.
\end{proof}

\begin{corollary}[Teichmuller dimension]
\label{thurston-ch13-teichmuller-dimension}
\dcref{13.3.7}
\group{surfaces}
\source{thurston:page-0336}
\notready
Let $O$ be a two-orbifold with $\chi(O)<0$, with $k$ elliptic points and $l$
corner reflectors.  Its Teichmuller space is homeomorphic to
\[
 \mathbb R^{-3\chi(X_O)+2k+l}.
\]
\end{corollary}
\begin{proof}
\sketch Decompose $O$ along closed geodesics and arcs perpendicular to its
boundary.  Lengths of the arcs, together with length and twist parameters on
the closed cuts, give Euclidean coordinates.  The formula holds for each
elementary piece; gluing circles preserves both sides of the formula, while
gluing arcs lowers both by one.
\end{proof}

\section{Fibrations and tetrahedral orbifolds}

\begin{definition}[Orbifold fibration]
\label{thurston-ch13-orbifold-fibration}
\dcref{13.4.1}
\group{fibrations}
\source{thurston:page-0336}
\notready
A fibration over an orbifold $O$ with generic fiber $F$ is locally a quotient
$(F\times\widetilde U)/\Gamma$, where the finite chart group acts on the base
and through structure-group automorphisms on $F$.
\end{definition}
\begin{proof}
\sketch Local products descend under the diagonal chart action, and changes of
orbifold charts induce the required structure-group transitions.
\end{proof}

\begin{proposition}[Geometry of tangent sphere orbifolds]
\label{thurston-ch13-unit-tangent-elliptic}
\dcref{13.4.2}
\group{fibrations}
\source{thurston:page-0337,thurston:page-0338}
\notready
For a two-orbifold $O$, the tangent construction has the following geometries:
\[
\begin{array}{c|c}
\text{geometry or type of }O&\text{geometry of its tangent orbifold}\\ \hline
\text{elliptic}&T_1(O)\text{ is elliptic of dimension three},\\
\text{Euclidean}&T_1(O)\text{ is Euclidean},\\
\text{bad}&TS(O)\text{ admits an elliptic structure}.
\end{array}
\]
\end{proposition}
\begin{proof}
\sketch In the elliptic case, identify $T_1(S^2)$ with $SO(3)$ and lift the
orbifold action to the double cover $S^3$.  In the Euclidean case use the
product geometry on $T_1(\mathbb E^2)=\mathbb E^2\times S^1$.  The tangent
sphere spaces of the bad orbifolds are lens spaces or reflection quotients of
lens spaces, and hence carry elliptic structures.
\end{proof}

\begin{definition}[Incompressible and Haken three-orbifolds]
\label{thurston-ch13-incompressible-haken}
\group{tetrahedra}
\source{thurston:page-0342}
\notready
Suppose that the underlying space $X_O$ of a three-orbifold $O$ is oriented.
A two-suborbifold $O'\subset O$ is incompressible if $\chi(O')\leq0$ and every
one-suborbifold of $O'$ which bounds an elliptic suborbifold of $O-O'$ also
bounds an elliptic suborbifold of $O'$.  The orbifold $O$ is Haken when it is
irreducible and contains an incompressible two-suborbifold.
\end{definition}

\begin{proposition}[Irreducibility criterion for tetrahedral orbifolds]
\label{thurston-ch13-tetrahedral-irreducible}
\dcref{13.5.1}
\group{tetrahedra}
\source{thurston:page-0342}
\notready
For an orbifold with underlying space $D^3$ and singular locus $\partial D^3$,
irreducibility is equivalent to the absence of an essential spherical
suborbifold; in the tetrahedral setting every spherical suborbifold bounds a
suborbifold ball.
\end{proposition}
\begin{proof}
\sketch Double across the boundary and use the finite spherical classification.
An embedded spherical section cuts the doubled tetrahedral model into balls,
which descend to the asserted ball-bounding property.
\end{proof}

\begin{proposition}[Haken criterion for reflection orbifolds]
\label{thurston-ch13-tetrahedral-haken}
\dcref{13.5.2}
\group{tetrahedra}
\source{thurston:page-0343,thurston:page-0344}
\uses{thurston-ch13-incompressible-haken}
\notready
An orbifold $O$ with $X_O=D^3$ and $\Sigma_O=\partial D^3$ is Haken if and only
if it is irreducible, is not the suspension of an elliptic two-orbifold, and
does not have the combinatorial type of a tetrahedron.
\end{proposition}
\begin{proof}
\sketch Encircle a face by a curve meeting only its incident singular edges.
The configurations with at most two such intersections force a suspension,
and an elliptic three-edge section forces tetrahedral type.  In every other
case the section is incompressible, or a compression produces an
incompressible triangular section.  Conversely, a section in a suspension or
a tetrahedral orbifold must revisit a face, or two faces sharing only one edge,
and is therefore compressible.
\end{proof}

\begin{theorem}[Simplex orbifold geometrization]
\label{thurston-ch13-simplex-geometrization}
\dcref{13.5.3}
\group{tetrahedra}
\source{thurston:page-0345,thurston:page-0346,thurston:page-0347}
\notready
Every $n$-orbifold having the combinatorial type of an $n$-simplex admits an
elliptic, Euclidean, or hyperbolic structure.
\end{theorem}
\begin{proof}
\sketch Form the matrix whose diagonal entries are $1$ and whose $ij$ entry is
$-\cos\alpha_{ij}$, where $\alpha_{ij}$ is the angle between two faces.  Every
principal submatrix obtained by deleting one row and column is positive
definite because it is the Gram matrix of an elliptic vertex link.  The full
matrix is therefore positive definite, positive semidefinite of rank $n$, or
has signature $(n,1)$, and these cases construct respectively the elliptic,
Euclidean, and hyperbolic simplex.
\end{proof}

\begin{theorem}[Hyperbolization of a simplex with deleted vertices]
\label{thurston-ch13-deleted-simplex-hyperbolization}
\dcref{13.5.4}
\group{tetrahedra}
\source{thurston:page-0346,thurston:page-0347}
\notready
Let $O$ be an $n$-orbifold with the combinatorial type of an $n$-simplex from
which some vertices have been deleted.  If every deleted vertex has a
Euclidean orbifold as its link and the Coxeter diagram of $O$ is connected,
then $O$ admits a complete finite-volume hyperbolic structure.
\end{theorem}
\begin{proof}
\sketch A deleted vertex gives a rank-$(n-1)$ principal submatrix of the face
Gram matrix.  In a basis adapted to its null vector, connectedness of the
Coxeter diagram makes the remaining coupling nonzero, so the full matrix has
signature $(n,1)$.  The other principal submatrices are positive definite or
rank-$(n-1)$ semidefinite, placing every vertex in hyperbolic space or on its
ideal boundary.  The resulting simplex is therefore complete and has finite
volume.
\end{proof}

\section{Convex polyhedra}

\begin{theorem}[Andreev's realization theorem]
\label{thurston-ch13-andreev}
\dcref{13.6.1}
\group{polyhedra}
\source{thurston:page-0348,thurston:page-0349,thurston:page-0350,thurston:page-0351}
\notready
Let $P$ be a combinatorial polyhedron with prescribed non-obtuse dihedral
angles.  A compact convex hyperbolic polyhedron realizing $P$ exists exactly
when the vertex, prismatic $3$-circuit, and prismatic $4$-circuit inequalities
hold, together with the stated strictness conditions; it is unique up to
hyperbolic isometry.
\end{theorem}
\begin{proof}
\sketch Truncate vertices, construct the dual circle pattern, and apply the
angle inequalities to rule out degenerating circuits.  Compactness gives a
limiting polyhedron, and convexity plus the dual pattern gives uniqueness.
\end{proof}

\begin{corollary}[Planar graph realization]
\label{thurston-ch13-planar-graph-realization}
\dcref{13.6.2,13.6.3}
\group{polyhedra}
\source{thurston:page-0348,thurston:page-0351}
\notready
For a planar graph satisfying the distinct-edge and circuit hypotheses of
Andreev's theorem, there is a convex hyperbolic polyhedron with that graph and
the prescribed admissible dihedral angles.  In particular every triangulation
of $S^2$ is realized by a convex polyhedron.
\end{corollary}
\begin{proof}
\sketch Apply Andreev to the graph after checking that its $3$- and $4$-circuits
are exactly the prohibited separating configurations; triangulations have no
additional prismatic obstruction.
\end{proof}

\begin{theorem}[Cusped orbifold realization]
\label{thurston-ch13-cusped-realization}
\dcref{13.6.4,13.6.5}
\group{polyhedra}
\source{thurston:page-0352,thurston:page-0353,thurston:page-0354}
\notready
The Andreev realization and uniqueness statements extend to ideal polygonal
orbifolds with torus or higher-genus ends, and to closed surfaces of negative
Euler characteristic with prescribed non-obtuse orbifold cell structure.
\end{theorem}
\begin{proof}
\sketch Replace ideal vertices by horospherical truncations and pass to the
limit.  The same circuit inequalities control compactness; the limiting
horospheres give the parabolic ends and uniqueness follows by duality.
\end{proof}

\begin{theorem}[Circle-pattern existence]
\label{thurston-ch13-circle-pattern}
\dcref{13.7.1}
\group{circle-patterns}
\source{thurston:page-0355,thurston:page-0356,thurston:page-0357,thurston:page-0358,thurston:page-0359,thurston:page-0360,thurston:page-0361,thurston:page-0362,thurston:page-0363,thurston:page-0364}
\notready
Let $S$ be closed with $\chi(S)\leq0$ and let $\tau$ be a cell division whose
face angles satisfy the non-obtuse compatibility conditions.  There is a
Euclidean or hyperbolic circle pattern on $S$ with incidence graph $\tau$ and
the prescribed intersection angles, unique up to similarity or isometry.
\end{theorem}
\begin{proof}
\sketch Assign radii to face circles.  Three circles with prescribed radii and
pairwise non-obtuse intersection angles exist uniquely in each geometry; the
resulting local angle equations have a strictly convex variational functional.
Its minimum supplies the pattern, while degeneration is excluded by the
Euler-characteristic hypotheses.  The Euclidean statement follows by taking
the zero-curvature limit.
\end{proof}

\begin{lemma}[Three-circle configuration]
\label{thurston-ch13-three-circle-configuration}
\dcref{13.7.2}
\group{circle-patterns}
\source{thurston:page-0356}
\notready
For positive radii $R_1,R_2,R_3$ and angles
$\theta_i\in[0,\pi/2]$, there is, in both hyperbolic and Euclidean geometry,
a configuration of three circles with those radii and pairwise intersection
angles, unique up to isometry.
\end{lemma}
\begin{proof}
\sketch The distance between centers is determined by the two radii and the
intersection angle through the cosine law.  The three distances determine a
triangle uniquely, and the non-obtuse bounds ensure its existence.
\end{proof}

\begin{lemma}[Circle-pattern convexity]
\label{thurston-ch13-circle-pattern-convexity}
\dcref{13.7.3}
\group{circle-patterns}
\source{thurston:page-0358,thurston:page-0359,thurston:page-0360}
\notready
The angle equations for the circle pattern are the gradient equations of a
strictly convex functional in logarithmic radii, modulo the Euclidean scaling
direction.  Consequently a solution is unique, and the variational argument
proves existence whenever the combinatorial angle inequalities hold.
\end{lemma}
\begin{proof}
\sketch Differentiate the three-circle cosine-law relation around each face.
The Hessian is a sum of positive semidefinite face contributions whose only
common kernel is constant scaling; the boundary behavior of the functional is
controlled by the circuit inequalities.
\end{proof}

\begin{corollary}[Euclidean and parabolic cases]
\label{thurston-ch13-circle-pattern-euclidean}
\dcref{13.7.4}
\group{circle-patterns}
\source{thurston:page-0364}
\notready
The realization theorems for the cusped orbifolds above remain valid when the
surface geometry is Euclidean or parabolic.
\end{corollary}
\begin{proof}
\sketch Apply the circle-pattern theorem in the zero-curvature limit and use
horospherical truncations at parabolic vertices.
\end{proof}

\section{Compactifications}

\begin{definition}[Hyperbolic structure with nodes]
\label{thurston-ch13-hyperbolic-nodes}
\group{compactification}
\source{thurston:page-0364,thurston:page-0365}
\notready
A hyperbolic structure with nodes on a two-orbifold is a complete finite-volume
hyperbolic structure on the complement of disjoint one-dimensional
suborbifolds, called nodes.  The topology up to isotopy is defined by
$e^\epsilon$-quasi-isometries outside fixed thin neighborhoods of the nodes.
\end{definition}
\begin{proof}
\sketch Removing shorter and shorter geodesics produces cusp neighborhoods;
quasi-isometry on the thick parts records convergence while disregarding the
length diverging collars.
\end{proof}

\begin{theorem}[Polygonal Teichmuller compactification]
\label{thurston-ch13-polygonal-compactification}
\dcref{13.8.1}
\group{compactification}
\source{thurston:page-0365,thurston:page-0366,thurston:page-0367,thurston:page-0368}
\notready
For an $n$-gonal orbifold $P$, the space $\mathcal N(P)$ of hyperbolic
structures with nodes is a closed disk $D^{n-3}$ whose interior is
$\mathcal T(P)$.  Its open cells are indexed by sets of pairwise disjoint nodes
up to isotopy.  Equivalently, its boundary cell structure is dual to the
polyhedral sphere of projective measured laminations on $P$.
\end{theorem}
\begin{proof}
\sketch Extend a maximal disjoint system of node arcs to geodesic arcs.  Their
lengths are bounded by area and converge, possibly to zero, giving compactness.
These lengths are local quadrant coordinates, and fixing the zero coordinates
gives Euclidean strata.  Convex-hull heights on a regular polygon identify
measured laminations with a piecewise-linear sphere; its dual convex body is
the required disk.
\end{proof}

\begin{definition}[Kleinian structures with nodes]
\label{thurston-ch13-kleinian-nodes}
\group{kleinian-compactification}
\source{thurston:page-0368,thurston:page-0369}
\notready
Let $O$ have underlying space $D^3$, singular locus in $\partial D^3$, and
polygonal $\partial\Sigma_O$.  A Kleinian structure with nodes is a Kleinian
structure off a disjoint union of nodes of three types: incompressible
one-suborbifolds of $\partial O$; incompressible Euler-zero two-suborbifolds
with nonempty boundary; and $P_{2k}\times\mathbb R$ suborbifolds with alternating
boundary and interior sides.  Node neighborhoods are horoball cusp
neighborhoods, and convergence is quasi-isometric on the thick convex-hull
part.
\end{definition}
\begin{proof}
\sketch These are exactly the possible thin cross-sections produced when a
geodesic or polygonal region becomes parabolic; removing horoball
neighborhoods leaves ordinary Kleinian quotient pieces.
\end{proof}

\begin{lemma}[Convex-hull volume bound]
\label{thurston-ch13-convex-hull-volume}
\dcref{13.9.2}
\group{kleinian-compactification}
\source{thurston:page-0369,thurston:page-0370}
\notready
There is a uniform upper bound for the volume of the convex hull of any
Kleinian structure with nodes on $O$.
\end{lemma}
\begin{proof}
\sketch The boundary bending lamination has uniformly boundedly many faces and
sides, hence uniformly bounded area.  The hyperbolic isoperimetric estimate
$\operatorname{vol}(S)\leq\tfrac12\operatorname{area}(\partial S)$ bounds the
convex-hull volume.
\end{proof}

\begin{theorem}[Kleinian deformation compactification]
\label{thurston-ch13-kleinian-compactification}
\dcref{13.9.1}
\group{kleinian-compactification}
\source{thurston:page-0369,thurston:page-0370,thurston:page-0371,thurston:page-0372,thurston:page-0373}
\notready
If $O$ is irreducible with incompressible boundary and admits one non-elementary
Kleinian structure, then its space $\mathcal N(O)$ of Kleinian structures with
nodes is compact.  The conformal-boundary map is a homeomorphism
\[
\mathcal N(O)\cong\mathcal N(\partial O).
\]
\end{theorem}
\begin{proof}
\sketch By the volume bound, degeneration can occur only through long edges of
the convex hull.  Thin cross-sections have nonnegative Euler characteristic and
are precisely cusp, annular, or boundary-compressible node pieces.  Removing
the resulting long spikes leaves bounded convex pieces, giving compactness.
The conformal-boundary map is continuous by thick-part quasi-isometries,
surjective by compactness and the boundary deformation theorem, and injective
by convex-hull rigidity; hence it is a homeomorphism.
\end{proof}
